% Part: first-order-logic
% Chapter: proof-systems
% Section: seqeunt-calculus

\documentclass[../../../include/open-logic-section]{subfiles}

\begin{document}

\iftag{FOL}
      {\olfileid{fol}{prf}{seq}}
      {\olfileid{pl}{prf}{seq}}

\olsection{సీక్వెంట్ కలనం}

అనేక !!{derivation} వ్యవస్థలు !!{sentence}s అమరికలతో పనిచేస్తే,
సీక్వెంట్ కలనం \emph{సీక్వెంట్లతో} పనిచేస్తుంది. సీక్వెంట్ అంటే
కింది రూపంలోని వ్యక్తీకరణ:
\[
!A_1, \dots, !A_m \Sequent !B_1, \dots, !B_n,
\]
అంటే సీక్వెంట్ సంకేతం~$\Sequent$తో వేరుచేసిన రెండు !!{sentence}s
క్రమాల జత. ఈ రెండు క్రమాల్లో దేనికైనా శూన్య క్రమం కావచ్చు.
\sourcecorrection{OLTEPRF-001}{సాధారణ సీక్వెంట్‌లోని రెండు క్రమాలకూ
ఆంగ్ల మూలం చివరి సూచికగా $m$నే ఇచ్చింది. వెంటనే వచ్చే నిర్వచనం
వాటిని స్వతంత్ర క్రమాలుగా తీసుకొని, దేనికైనా శూన్య క్రమం కావచ్చని
చెబుతుంది. అందువల్ల కుడి క్రమం చివరి సూచికను $n$గా సరిచేశాం.}
సీక్వెంట్ కలనంలోని ఒక !!^a{derivation} సీక్వెంట్ల వృక్షం. అందులో
అన్నిటికంటే పైనున్న సీక్వెంట్లు ప్రత్యేక రూపంలో ఉంటాయి (వాటిని
“ప్రారంభ సీక్వెంట్లు” లేదా “స్వీకృతాలు” అంటారు); మిగతా ప్రతి
సీక్వెంట్ తనకు వెంటనే పైనున్న సీక్వెంట్ల నుంచి ఏదో ఒక నిగమన నియమం
ద్వారా వస్తుంది. నిగమన నియమాలు సీక్వెంట్లలోని !!{sentence}sను
(ఎడమవైపున లేదా కుడివైపున చేర్చడం, తొలగించడం, పునర్వ్యవస్థీకరించడం
ద్వారా) మార్చవచ్చు, లేదా నియమపు నిష్కర్షలో ఒక సంక్లిష్ట
!!{formula}ను ప్రవేశపెట్టవచ్చు. ఉదాహరణకు, $\LeftR{\land}$ నియమం
$!A, \Gamma \Sequent \Delta$ నుంచి $!A \land !B, \Gamma \Sequent
\Delta$ను నిగమించనిస్తుంది; $\RightR{\lif}$ నియమం ఏ $\Gamma$,
$\Delta$, $!A$, $!B$లకైనా $!A, \Gamma \Sequent \Delta, !B$ నుంచి
$\Gamma \Sequent \Delta, !A \lif !B$ను నిగమించనిస్తుంది.
(ప్రత్యేకంగా, $\Gamma$, $\Delta$లు శూన్య క్రమాలు కావచ్చు.)

సీక్వెంట్ కలనం ఆధారంగా $\Proves$ సంబంధాన్ని ఇలా నిర్వచిస్తాం:
$\Gamma_0$లోని ప్రతి $!A$, $\Gamma$లో ఉండి, మూలం వద్ద
$\Gamma_0 \Sequent !A$ సీక్వెంట్ గల ఒక !!{derivation} ఉండేలా ఏదో
అలాంటి క్రమం ఉంటే, అప్పుడు మాత్రమే $\Gamma \Proves !A$.
$\Sequent !A$ సీక్వెంట్‌కు ఒక !!a{derivation} ఉంటే $!A$ సీక్వెంట్
కలనంలో ఒక సిద్ధాంతం. ఉదాహరణకు, $\Proves (!A \land !B) \lif !A$ అని
చూపే ఒక !!a{derivation} ఇది:
\begin{prooftree}
  \Axiom$!A \fCenter !A$
  \RightLabel{\LeftR{\land}}
  \UnaryInf$!A \land !B \fCenter !A$
  \RightLabel{\RightR{\lif}}
  \UnaryInf$\fCenter (!A \land !B) \lif !A$
\end{prooftree}

$\Gamma_0$లోని ప్రతి $!A \in \Gamma_0$, $\Gamma$లో ఉండి,
$\Gamma_0 \Sequent$కు (సీక్వెంట్ కుడివైపు శూన్యంగా ఉండే) ఒక
!!a{derivation} ఉంటే, $\Gamma$ సీక్వెంట్ కలనంలో వైరుధ్యభరితమైనది.
\RightR{\Weakening} నియమాన్ని ఉపయోగించి, వైరుధ్యభరిత సమితి నుంచి
ఏ !!{sentence}నైనా !!{derive} చేయవచ్చు.

సీక్వెంట్ కలనాన్ని 1930లలో గెర్హార్డ్ గెంట్సెన్ కనుగొన్నాడు. దాని
క్రమబద్ధమైన, సౌష్ఠవమైన రూపకల్పన వల్ల !!{derivation}s సిద్ధాంతాన్ని
అభివృద్ధి చేయడానికి అది ఎంతో ఉపయోగకరమైన వ్యవస్థ. సీక్వెంట్ కలనంలో
!!{derivation}sను కనుగొనడం సాపేక్షంగా సులభం; కానీ ఈ
!!{derivation}sను చదవడం
తరచుగా కష్టం, నిరూపణలతో వాటి సంబంధం కొన్నిసార్లు సులభంగా
కనిపించదు. అయినప్పటికీ, !!{derivation} వ్యవస్థలకు ఇది ఎంతో సొగసైన
విధానంగా నిరూపితమైంది; అనేక తర్కవ్యవస్థలకు సీక్వెంట్ కలన
వ్యవస్థలు ఉన్నాయి.

\end{document}
